% ===== §10 Directed Fissures (the Anti-Envoi) =====
\section{Directed Fissures}
\label{sec:fissures}

\noindent
A closing section is, by convention, where an argument gathers itself: collects its results, resolves its tensions, comes to rest. This one does the opposite, and the difference is a matter of the paper's place in the series rather than of style. Were this the final paper, the figure of water would be made to bring the long argument home, to settle it in the Way; ``close to the Way'' (ch.~8) would be the resting point. But this is a paper near the middle, and water was chosen not because it concludes but because it is a source --- a figure rich enough that the pressure of pursuing it past the edge of what this paper can hold opens, at four definite places, the seeds of papers still to come. We name the four fissures rather than seal them. Each is a point where the water-virtue, pressed further, runs out of what the present argument establishes and points toward what a further argument would have to establish. Water here is the spring, not the sea.

\subsection{Watershed justice}
\label{sub:fissure-basin}

The first fissure is the one §8 opened and declined to cross. The water-virtue has been established between two; whether it governs the many --- the family, the community, the relational commons --- is a substantive question with its own difficulties, not a corollary. The non-substitutable flesh that grounds intimate particularity may have no analogue at the scale of types and roles; the rigid lower bound may not translate where the self at risk of being drained is a class rather than a person; and the holonomy criterion is, in the formal companion, sensitive to the number of subjects and the emergence of their shared background in ways that bear directly on whether a two-body criterion governs an n-body commons \citep{wan2026braid,paperxv}. The materials for crossing exist --- Ostrom's governance of the literal watershed \citep{ostrom1990}, the land ethic's enlargement of the unit of benefit \citep{leopold1949} --- but the crossing is the work of another paper. A watershed justice: the water-virtue scaled from the coupling to the basin, with the distortions of scaling faced rather than assumed away.

\subsection{The temporal dynamics of water}
\label{sub:fissure-time}

The second fissure is in the dimension this paper has treated only at rest. The water-virtue has been given largely as a structure --- a geometry of flow, a criterion on circulation --- but water is above all a thing in time, and a relation's water-good is not a fixed state but a process of attunement that drifts, falls out of phase, and must be re-tuned. The resonance of two meaning-worlds \citep{paperxix} is a phase relation, and phase relations have a temporal life: they synchronise, they slip, they require continual re-coupling to hold. The formal companion is explicit that the dynamical layer of its own construction --- how value flows in time, at what rate the structure forms and dissolves, how a background of relation emerges --- is left open \citep{wan2026braid}. A paper on the temporal water-virtue would take up exactly this: attunement and mis-attunement as the lived time of intimacy, the work of re-tuning a resonance that no relation holds without effort, the difference between a phase that is being continually renewed and one that has frozen. Non-fullness (\xref{sec:nonfullness}) is the static shadow of this temporal virtue; its full form is dynamical and unwritten.

\subsection{Water and vessel}
\label{sub:fissure-vessel}

The third fissure concerns the relation between the formless flow and the form that shapes it. Throughout this paper the vessel has appeared as the thing whose shape water takes; but the vessel is also what gives water a determinate course, and the relation between a flow and the form that channels it is the relation between an intimacy and the institution --- the vow, the contract, the structured commitment --- that gives it shape without freezing it. This series has treated the contract as structured flow on a manifold and the non-binding vow as the form of a commitment that holds without an external enforcer \citep{paperxiii,paperxix}. A paper on water and vessel would ask how form may give the formless an edge without solidifying it: how a commitment can shape a flow toward the returning geometry (\xref{sub:formal-power}) without hardening into the fixed channel that is solidification's bad form (\xref{sub:formal-solidification}). The vessel that channels without freezing, that gives course without capture, is the institutional question the water-virtue raises and this paper leaves open.

\subsection{The subject of the highest good}
\label{sub:fissure-subject}

The fourth fissure is the one the formal incompleteness opened (\xref{sub:formal-incompleteness}), and it is the deepest. If the water-good cannot be reduced to a procedure --- if judgment in the flesh is irreducible --- then there must be a subject who judges, and the nature of that subject is not settled by anything in this paper. The danger is sharp and was named in the survey: ``taking the vessel's shape'' may be the water-virtue or the self-erasure of a subject who has given way on its desire \citep{lacan2006}, and the difference between the one who is like water and the one who has been dissolved into water is a difference in the subject, not in the conduct. A paper on the subject of the highest good would ask who the one who is like water without being effaced must be --- what kind of self can yield without dissolving, descend without submitting, empty without disappearing; how the relational subject of this series' founding commitments \citep{wan2024grb} holds a desire of its own while taking, genuinely, the shape of the other. The rigid lower bound (\xref{sub:dialectic-bound}) is the criterion such a subject must satisfy; what such a subject \emph{is} --- the one who can be water and remain --- is the question with which the water-virtue finally hands itself on. It is the spring's last opening, and the one toward which, perhaps, the rest of the series runs.
